\documentclass[11pt]{article}

\usepackage[margin=1in]{geometry}
\usepackage[T1]{fontenc}
\usepackage[utf8]{inputenc}
\usepackage{amsmath,amssymb,amsthm,mathtools}
\usepackage{bm}
\usepackage{booktabs}
\usepackage{array}
\usepackage{longtable}
\usepackage{tabularx}
\usepackage{multirow}
\usepackage{xcolor}
\usepackage{enumitem}
\usepackage{hyperref}
\usepackage{float}
\usepackage{caption}
\usepackage{titlesec}
\usepackage{listings}

\hypersetup{
  colorlinks=true,
  linkcolor=blue,
  urlcolor=blue,
  citecolor=blue
}

\titleformat{\section}{\large\bfseries}{\thesection}{0.6em}{}
\titleformat{\subsection}{\normalsize\bfseries}{\thesubsection}{0.6em}{}
\titleformat{\subsubsection}{\normalsize\bfseries}{\thesubsubsection}{0.6em}{}

\setlist[itemize]{topsep=3pt,itemsep=2pt,leftmargin=18pt}
\setlist[enumerate]{topsep=3pt,itemsep=2pt,leftmargin=20pt}

\lstset{
  basicstyle=\ttfamily\small,
  frame=single,
  breaklines=true,
  columns=fullflexible
}

\definecolor{decisionblue}{RGB}{41,98,255}
\definecolor{notegreen}{RGB}{0,150,80}
\definecolor{warnorange}{RGB}{230,120,0}

\newcommand{\decision}[1]{\textcolor{decisionblue}{\textbf{Decision:} #1}}
\newcommand{\note}[1]{\textcolor{notegreen}{\textbf{Note:} #1}}
\newcommand{\warn}[1]{\textcolor{warnorange}{\textbf{Warning:} #1}}
\newcommand{\openq}[1]{\textcolor{warnorange}{\textbf{Open:} #1}}

\title{%
  Detailed ML Architecture Specification\\[4pt]
  \large Real-Time Adaptive Music Recommender --- v5.0\\[2pt]
  \normalsize Team Spoty Boys}
\author{}
\date{\today}

\begin{document}
\maketitle
\tableofcontents
\newpage

%% ===================================================================
\section{Overview of Changes from v4.0}
%% ===================================================================

Version 5.0 replaces the PANNs audio embedding pipeline with
\textbf{Item2Vec}. The primary motivation is engineering feasibility:
the PANNs approach requires downloading approximately 4.5 million
30-second MP3 preview clips from Deezer, a process estimated to take
$\sim$10 days of continuous operation with a hit rate of only $\sim$30\%.
Item2Vec eliminates all external downloads and trains entirely on
the session interaction data already on disk.

\begin{table}[H]
\centering
\caption{Key differences between v4.0 (PANNs) and v5.0 (Item2Vec).}
\label{tab:version_diff}
\begin{tabular}{lll}
\toprule
\textbf{Property} & \textbf{v4.0 (PANNs)} & \textbf{v5.0 (Item2Vec)} \\
\midrule
Embedding source & Acoustic content (audio waveform) & User behaviour (session co-occurrence) \\
External data required & $\sim$4.5M MP3 downloads & None \\
Preprocessing time & $\sim$10 days (download) + 4--6h (GPU) & $\sim$1 hour (CPU only) \\
Catalog coverage & $\sim$30\% of session tracks & $\sim$80--90\% of session interactions \\
Embedding dimension & 2048 $\to$ 128 (TruncatedSVD) & 128 (direct) \\
Cold-item handling & Any track with an MP3 clip & Tracks with $\ge 5$ session appearances \\
GPU required & Yes (PANNs forward pass) & No \\
\midrule
ALS branch & Removed (v4.0 decision) & Not added back \\
Tag features & Removed (v4.0 decision) & Not added back \\
Scalar features & Removed (v4.0 decision) & Not added back \\
\bottomrule
\end{tabular}
\end{table}

All other architectural decisions from v4.0 --- the GRU ranker, the
three-branch retriever (co-occurrence, Pref~NN, popularity), the
multi-centroid user preference vector, the joint BCE\,+\,BPR training
objective, and the policy re-ranker --- are carried forward unchanged.
The document below is self-contained; v4.0 need not be consulted.

%% ===================================================================
\section{Raw Data Inventory}
%% ===================================================================

All raw files are located at
\texttt{data/raw/content/30music\_parsed/}.
Table~\ref{tab:raw_inventory} lists every file, its row count, and
its role in the pipeline.

\begin{table}[H]
\centering
\caption{Raw file inventory.}
\label{tab:raw_inventory}
\begin{tabular}{lrl}
\toprule
\textbf{File} & \textbf{Rows} & \textbf{Role} \\
\midrule
\texttt{tracks.csv}           & 5,675,143  & Track catalog; deduplicated to obtain unique tracks \\
\texttt{events.csv}           & ---        & Raw play log; used to compute per-track play counts \\
\texttt{persons.csv}          & 595,049    & Artist reference; not used in model \\
\texttt{albums.csv}           & 217,337    & Not used \\
\texttt{tags.csv}             & 276,618    & Not used (70.5\% of tracks have no tags) \\
\texttt{users.csv}            & 45,167     & User reference; only \texttt{user\_id} needed \\
\texttt{love.csv}             & 4,106,341  & Explicit positive signal ($3\times$ weight in user centroids) \\
\texttt{playlist\_meta.csv}   & 57,561     & Playlist ownership \\
\texttt{playlist\_tracks.csv} & 1,603,040  & Playlist-track membership; co-occurrence construction \\
\texttt{session\_meta.csv}    & 2,764,474  & Session ownership \\
\texttt{session\_tracks.csv}  & 31,351,945 & Core interaction data; Item2Vec corpus and model training \\
\bottomrule
\end{tabular}
\end{table}

\subsection{Key Interaction Statistics from Data Inspection}

\begin{table}[H]
\centering
\caption{Session interaction statistics (from \texttt{data\_inspection\_report}).}
\begin{tabular}{lr}
\toprule
\textbf{Quantity} & \textbf{Value} \\
\midrule
Total sessions & 2,764,474 \\
Sessions with length $\ge 2$ & 2,249,528 \\
Adjacent positive/neutral pairs & 21,908,813 \\
Unique track IDs in sessions & 3,324,298 \\
Session median length & 6 \\
Session mean length & 11.3 \\
Users with $\ge 50$ interactions & 90.9\% \\
\midrule
ALS positive-only interaction density & 0.0098\% (595 GB float32) \\
\bottomrule
\end{tabular}
\end{table}

\begin{table}[H]
\centering
\caption{Label distribution in \texttt{session\_tracks.csv}.}
\begin{tabular}{lrr}
\toprule
\textbf{Label} & \textbf{Count} & \textbf{Fraction} \\
\midrule
\texttt{positive} & $\sim$23,188,858 & 73.9\% \\
\texttt{unknown}  & $\sim$4,734,144  & 15.1\% \\
\texttt{neutral}  & $\sim$2,194,636  & 7.0\%  \\
\texttt{skip}     & $\sim$1,254,078  & 4.0\%  \\
\midrule
Total & 31,351,945 & 100\% \\
\bottomrule
\end{tabular}
\end{table}

\subsection{Why ALS and Tag Features Are Excluded}

\paragraph{ALS.}
The positive-only interaction matrix spans 44,807 users $\times$ 3,324,298 tracks
at a density of $0.0098\%$. In this extreme sparsity regime, ALS latent factors
for long-tail tracks are numerically unreliable (the gradient signal is
dominated by a tiny fraction of popular items). The full positive-and-love
matrix has essentially the same density ($0.0098\%$). ALS is therefore excluded.

\paragraph{Tag features.}
70.5\% of tracks in the catalog have no tag annotations whatsoever.
A TF-IDF/SVD encoding over such a sparse, skewed vocabulary would produce
near-zero vectors for the majority of tracks and is not expected to add
useful signal. Tag features are excluded.

%% ===================================================================
\section{Data Preprocessing Pipeline}
%% ===================================================================

\subsection{Pipeline Overview}

The preprocessing pipeline has seven stages executed in the order shown
in Figure~\ref{fig:preproc_overview}. Stages~1--3 build the track catalog;
Stages~4--6 train the Item2Vec model and define the \emph{item2vec catalog}
(the set of tracks for which embeddings exist); Stage~7 filters all
interaction tables to this catalog.

\begin{figure}[H]
\centering
\small
\begin{tabular}{cl}
\textbf{Stage} & \textbf{Description} \\
\midrule
1 & \texttt{tracks.csv} $\to$ unique track catalog (\texttt{tracks\_unique.csv}) \\
2 & \texttt{events.csv} $\to$ per-track event play count \\
3 & Merge catalog and play counts $\to$ \texttt{catalog.csv} \\
4 (A) & \texttt{session\_tracks.csv} $\to$ Item2Vec training corpus \\
5 (B) & Train Item2Vec model $\to$ embeddings + \texttt{item2vec\_catalog.csv} \\
6 (C) & Validate embeddings \\
7 (D) & Filter all interaction tables to \texttt{item2vec\_catalog} \\
\end{tabular}
\caption{Preprocessing stage order. Stages~4--7 correspond to sub-modules
\texttt{stage\_a} through \texttt{stage\_d} in
\texttt{src/data\_pre\_process/item2vec/}.}
\label{fig:preproc_overview}
\end{figure}

\decision{Item2Vec is trained on \textbf{all available sessions}, not only
the training split. Item2Vec is an unsupervised model with no behavioural
labels; it learns only the co-occurrence graph structure. Using all sessions
maximises vocabulary coverage and embedding quality. Unlike C3 (which is trained
with supervision and must not observe val/test labels), Item2Vec has no
information-leakage risk from this choice.}

\subsection{Stage 1 --- Unique Track Catalog}

\textbf{Problem.}
\texttt{tracks.csv} contains 5,675,143 rows but is not guaranteed to be
unique per \texttt{track\_id}. Duplicate entries arise from the crawling process.

\begin{table}[H]
\centering
\caption{Column selection for \texttt{tracks.csv}.}
\begin{tabular}{lll}
\toprule
\textbf{Column} & \textbf{Keep?} & \textbf{Reason} \\
\midrule
\texttt{track\_id}    & Yes & Primary key throughout the pipeline \\
\texttt{artist\_hint} & Yes & Human-readable name for debugging and display \\
\texttt{title}        & Yes & Human-readable name for debugging and display \\
\texttt{mbid}         & No  & 78.9\% null; not used \\
\texttt{duration}     & No  & 48.0\% are $-1$; not used \\
\texttt{playcount}    & No  & 22.5\% null and negative values observed; superseded by Stage~3 \\
\texttt{album\_ids}   & No  & Not used \\
\texttt{artist\_ids}  & No  & Not used \\
\texttt{tag\_ids}     & No  & Tag features excluded; 70.5\% tracks have no tags \\
\bottomrule
\end{tabular}
\end{table}

\begin{lstlisting}
import pandas as pd

tracks = pd.read_csv("tracks.csv",
                     usecols=["track_id", "artist_hint", "title"])
tracks = tracks.drop_duplicates(subset="track_id", keep="first")
tracks.to_csv("processed/tracks_unique.csv", index=False)
\end{lstlisting}

\textbf{Output:} \texttt{processed/tracks\_unique.csv} ---
one row per unique \texttt{track\_id}, columns
\texttt{track\_id}, \texttt{artist\_hint}, \texttt{title}.

\subsection{Stage 2 --- Per-Track Event Play Count}

\textbf{Motivation.}
The \texttt{playcount} column in \texttt{tracks.csv} has 22.5\% nulls and
observed negative values. \texttt{events.csv} is a direct play log; counting
rows per \texttt{track\_id} gives a clean, independently derived play count.

\begin{lstlisting}
events = pd.read_csv("events.csv", usecols=["track_id"])
event_counts = (events.groupby("track_id")
                       .size()
                       .reset_index(name="event_count"))
event_counts.to_csv("processed/event_playcount.csv", index=False)
\end{lstlisting}

\textbf{Output:} \texttt{processed/event\_playcount.csv} ---
columns \texttt{track\_id}, \texttt{event\_count}.

\subsection{Stage 3 --- Merged Track Catalog}

Join Stage~1 and Stage~2 on \texttt{track\_id}. Tracks absent from
\texttt{events.csv} receive \texttt{final\_playcount} = 0.

\begin{table}[H]
\centering
\caption{Derived fields added to the catalog.}
\begin{tabular}{lll}
\toprule
\textbf{Field} & \textbf{Formula} & \textbf{Used for} \\
\midrule
\texttt{pop\_score}          & $\log(1 + \texttt{final\_playcount})$ &
  Popularity branch; C4 fallback \\
\texttt{neg\_sample\_weight} & $\texttt{final\_playcount}^{0.75}$, normalised to $\sum = 1$ &
  C3 random negative sampling \\
\bottomrule
\end{tabular}
\end{table}

\begin{lstlisting}
tracks = pd.read_csv("processed/tracks_unique.csv")
event_counts = pd.read_csv("processed/event_playcount.csv")

catalog = tracks.merge(event_counts, on="track_id", how="left")
catalog["final_playcount"] = catalog["event_count"].fillna(0).astype(int)
catalog = catalog.drop(columns=["event_count"])

import numpy as np
catalog["pop_score"] = np.log1p(catalog["final_playcount"])
weights = catalog["final_playcount"].clip(lower=0).pow(0.75)
catalog["neg_sample_weight"] = weights / weights.sum()

catalog.to_csv("processed/catalog.csv", index=False)
\end{lstlisting}

\textbf{Output:} \texttt{processed/catalog.csv} ---
columns \texttt{track\_id}, \texttt{artist\_hint}, \texttt{title},
\texttt{final\_playcount}, \texttt{pop\_score}, \texttt{neg\_sample\_weight}.

\subsection{Stage 4 (A) --- Item2Vec Training Corpus}

\textbf{Purpose.}
Transform \texttt{session\_tracks.csv} into a list of ordered track-ID sequences
suitable for Word2Vec training. Each session becomes one sentence;
each \texttt{track\_id} becomes one word.

\textbf{Label filtering rationale.}
\begin{itemize}
  \item \textbf{Positive} ($73.9\%$): completed or substantially completed
        listening. Strong co-occurrence signal.
  \item \textbf{Neutral} ($7.0\%$): partial listening but not a skip.
        Retained as weaker co-occurrence evidence.
  \item \textbf{Skip} ($4.0\%$): explicit negative feedback. Including skip
        events would train the model that skip-pairs are semantically related,
        which is the opposite of the desired signal. Excluded.
  \item \textbf{Unknown} ($15.1\%$): playratio is null; no interpretable
        signal. Excluded.
\end{itemize}

\begin{table}[H]
\centering
\caption{Column selection for \texttt{session\_tracks.csv} (corpus stage).}
\begin{tabular}{lll}
\toprule
\textbf{Column} & \textbf{Keep?} & \textbf{Reason} \\
\midrule
\texttt{session\_id} & Yes & Group key \\
\texttt{position}    & Yes & Defines playback order within session \\
\texttt{track\_id}   & Yes & The vocabulary token \\
\texttt{label}       & Yes & Used to filter rows; not passed to the model \\
\texttt{user\_id}    & No  & Not needed for corpus construction \\
\texttt{playratio}   & No  & Not needed \\
\texttt{playstart}   & No  & Not needed \\
\texttt{playtime}    & No  & Not needed \\
\texttt{action}      & No  & Not needed \\
\bottomrule
\end{tabular}
\end{table}

\begin{lstlisting}
import pandas as pd

CHUNK_SIZE = 500_000
kept_labels = {"positive", "neutral"}
sessions = {}  # session_id -> sorted list of track_ids

for chunk in pd.read_csv(
        "session_tracks.csv",
        usecols=["session_id", "position", "track_id", "label"],
        chunksize=CHUNK_SIZE):
    chunk = chunk[chunk["label"].isin(kept_labels)]
    chunk = chunk.dropna(subset=["track_id"])
    chunk["track_id"] = chunk["track_id"].astype(int)
    for sid, grp in chunk.groupby("session_id"):
        grp_sorted = grp.sort_values("position")["track_id"].tolist()
        if sid in sessions:
            sessions[sid].extend(grp_sorted)
        else:
            sessions[sid] = grp_sorted

# Build dataframe and drop short sessions
rows = [(sid, ids, len(ids))
        for sid, ids in sessions.items() if len(ids) >= 2]
corpus = pd.DataFrame(rows, columns=["session_id", "track_ids", "length"])
corpus.to_parquet("processed/item2vec_corpus.parquet", index=False)

print(f"Sessions: {len(corpus):,}")
print(f"Total tokens: {corpus['length'].sum():,}")
print(f"Unique tracks: {len({t for ids in corpus['track_ids'] for t in ids}):,}")
\end{lstlisting}

\textbf{Expected output statistics.}
\begin{itemize}
  \item $\approx$2.1\,M sessions (from $\approx$2.25\,M usable sessions,
        after removing those that become length~$< 2$ post-filtering).
  \item $\approx$24\,M tokens.
  \item Vocabulary before \texttt{min\_count} filter: $\approx$3.3\,M unique tracks.
\end{itemize}

\textbf{Output:} \texttt{processed/item2vec\_corpus.parquet} ---
columns \texttt{session\_id}, \texttt{track\_ids} (list of int), \texttt{length}.

\subsection{Stage 5 (B) --- Item2Vec Model Training}

\textbf{Algorithm.}
Item2Vec treats each session sequence as a sentence and each
\texttt{track\_id} as a word, then applies the Word2Vec Skip-gram
objective. The Skip-gram model maximises the log-probability of
observing context tracks within a sliding window given a centre track:
\[
  \mathcal{L}_{\text{Skip-gram}} = \sum_{\text{session}} \sum_{t}
    \sum_{c \in W(t)} \log P(i_c \mid i_t),
\]
where $W(t)$ is the window of $\pm$\texttt{window} positions around
position $t$, and $P(i_c \mid i_t)$ is approximated via negative sampling.

\textbf{Why Skip-gram over CBOW.}
Skip-gram predicts context items from a centre item; it performs
better on rare items because each training example focuses on one
item at a time. In our long-tail catalog, many tracks appear
infrequently, making Skip-gram the preferred variant.

\begin{table}[H]
\centering
\caption{Item2Vec hyperparameters and their justification.}
\label{tab:i2v_hyperparams}
\begin{tabular}{lll}
\toprule
\textbf{Hyperparameter} & \textbf{Value} & \textbf{Justification} \\
\midrule
\texttt{vector\_size} & 128 &
  Compact; direct input to FAISS and C3 (no reduction needed) \\
\texttt{window} & 10 &
  Covers session median (6) and mean (11.3) in both directions \\
\texttt{min\_count} & 5 &
  Removes tracks with insufficient co-occurrence evidence;
  defines \texttt{item2vec\_catalog} \\
\texttt{sg} & 1 &
  Skip-gram; superior for rare items in long-tail distributions \\
\texttt{negative} & 15 &
  Noise-contrastive estimation; sampling $\propto$
  \texttt{freq}$^{0.75}$ (from \texttt{neg\_sample\_weight}) \\
\texttt{epochs} & 10 & Standard for music co-occurrence data \\
\texttt{workers} & 8 & Parallel training threads \\
\texttt{seed} & 42 & Reproducibility \\
\bottomrule
\end{tabular}
\end{table}

\begin{lstlisting}
import pandas as pd
import numpy as np
import json
from gensim.models import Word2Vec

# Load corpus
corpus_df = pd.read_parquet("processed/item2vec_corpus.parquet")
sentences = [list(map(str, ids)) for ids in corpus_df["track_ids"]]

# Train
model = Word2Vec(
    sentences=sentences,
    vector_size=128,
    window=10,
    min_count=5,
    sg=1,           # Skip-gram
    negative=15,
    epochs=10,
    workers=8,
    seed=42,
)
model.wv.save("processed/item2vec_model.bin")

# Extract embedding matrix in a deterministic row order
vocab_track_ids = [int(w) for w in model.wv.index_to_key]
emb_matrix = np.array([model.wv[str(t)] for t in vocab_track_ids],
                      dtype="float32")
np.save("processed/item2vec_128d.npy", emb_matrix)

# Row index mapping: track_id -> row in emb_matrix
track_id_to_row = {str(t): i for i, t in enumerate(vocab_track_ids)}
json.dump(track_id_to_row,
          open("processed/item2vec_track_to_row.json", "w"))

# Item2Vec catalog: join with catalog.csv for metadata
import pandas as pd
catalog = pd.read_csv("processed/catalog.csv")
i2v_catalog = catalog[catalog["track_id"].isin(set(vocab_track_ids))].copy()
i2v_catalog["i2v_row"] = i2v_catalog["track_id"].map(
    {t: i for i, t in enumerate(vocab_track_ids)})
i2v_catalog.to_csv("processed/item2vec_catalog.csv", index=False)

print(f"Vocabulary size: {len(vocab_track_ids):,}")
print(f"Item2Vec catalog size: {len(i2v_catalog):,}")
\end{lstlisting}

\textbf{Expected output statistics.}
\begin{itemize}
  \item Vocabulary size after \texttt{min\_count=5}: $\approx$400\,000--600\,000 tracks.
  \item Training time: 30--60 minutes on CPU with 8 workers.
  \item Coverage of session interactions: $\approx$80--90\% (popular tracks
        account for most plays; the long tail below \texttt{min\_count=5}
        represents a small fraction of interactions).
\end{itemize}

\textbf{Outputs:}
\begin{itemize}
  \item \texttt{processed/item2vec\_model.bin}: saved \texttt{KeyedVectors}.
  \item \texttt{processed/item2vec\_128d.npy}: embedding matrix,
        shape $(\text{vocab\_size}, 128)$.
  \item \texttt{processed/item2vec\_track\_to\_row.json}:
        \texttt{track\_id} $\to$ row index mapping.
  \item \texttt{processed/item2vec\_catalog.csv}:
        metadata for all embedded tracks.
\end{itemize}

\subsection{Stage 6 (C) --- Embedding Validation}

This stage performs automated checks to catch training failures before
downstream components are built.

\begin{enumerate}
  \item \textbf{Same-artist similarity check.}
        Sample 200 same-artist track pairs. Verify that their mean cosine
        similarity is higher than the mean cosine similarity of 200
        random track pairs. Failure indicates that the model has not
        learned musically coherent structure.
  \item \textbf{Numerical health check.}
        Assert: no NaN or infinite values in \texttt{item2vec\_128d.npy};
        embedding mean $\approx 0$ (Skip-gram produces approximately
        zero-centred vectors); standard deviation $\in (0.1, 2.0)$.
  \item \textbf{Nearest-neighbour spot check.}
        For 20 randomly sampled tracks, retrieve the top-5 nearest
        neighbours by cosine similarity. Inspect for plausibility
        (same artist, similar genre or era). Log results to
        \texttt{logs/item2vec\_validate.log}.
\end{enumerate}

\begin{lstlisting}
import numpy as np
from gensim.models import KeyedVectors

wv = KeyedVectors.load("processed/item2vec_model.bin")

# 1. Same-artist similarity
catalog = pd.read_csv("processed/item2vec_catalog.csv")
# ... (group by artist_hint, sample pairs, compute cosine similarity)

# 2. Numerical health
emb = np.load("processed/item2vec_128d.npy")
assert not np.isnan(emb).any(), "NaN values found in embeddings"
assert not np.isinf(emb).any(), "Inf values found in embeddings"
print(f"Mean: {emb.mean():.4f}, Std: {emb.std():.4f}")

# 3. Nearest-neighbour spot check
sample_ids = catalog["track_id"].sample(20).astype(str).tolist()
for tid in sample_ids:
    if tid in wv:
        neighbours = wv.most_similar(tid, topn=5)
        print(f"Track {tid}: {neighbours}")
\end{lstlisting}

\textbf{Output:} \texttt{logs/item2vec\_validate.log}.

\subsection{Stage 7 (D) --- Filter Interaction Tables to Item2Vec Catalog}

\decision{All interaction tables are filtered to retain only tracks present
in \texttt{item2vec\_catalog.csv} (i.e., tracks for which an Item2Vec
embedding exists). Tracks below the \texttt{min\_count=5} threshold are
removed from all tables. Sessions and playlists that become length
$\le 1$ after filtering are dropped entirely.}

\begin{lstlisting}
i2v_catalog = pd.read_csv("processed/item2vec_catalog.csv")
i2v_ids = set(i2v_catalog["track_id"])

def filter_and_drop_short(df, group_col, min_len=2):
    df = df[df["track_id"].isin(i2v_ids)].copy()
    counts = df.groupby(group_col)["track_id"].transform("count")
    return df[counts >= min_len].reset_index(drop=True)
\end{lstlisting}

\subsubsection{\texttt{session\_tracks.csv}}

\begin{table}[H]
\centering
\caption{Column selection for \texttt{session\_tracks.csv} (filtering stage).}
\begin{tabular}{lll}
\toprule
\textbf{Column} & \textbf{Keep?} & \textbf{Reason} \\
\midrule
\texttt{session\_id} & Yes & Group key \\
\texttt{user\_id}    & Yes & User identity \\
\texttt{position}    & Yes & Defines sequence order \\
\texttt{track\_id}   & Yes & Core entity \\
\texttt{playratio}   & Yes & Used to derive label and validate \\
\texttt{label}       & Yes & Training signal for C3 \\
\texttt{playstart}   & No  & Relative time within session; unusable as absolute \\
\texttt{playtime}    & No  & Not used \\
\texttt{action}      & No  & 9.3\% null; superseded by \texttt{label} \\
\bottomrule
\end{tabular}
\end{table}

\textbf{Additional cleaning steps.}
\begin{enumerate}
  \item Remove rows where \texttt{track\_id} $\notin$ \texttt{item2vec\_catalog}.
  \item Remove rows with \texttt{label = unknown} (playratio is null; no training signal).
  \item Cap \texttt{playratio} at 5.0 (210,483 rows exceed this; likely looped playback).
  \item Drop sessions where the remaining track count $< 2$.
  \item Re-sort by \texttt{(session\_id, position)}.
\end{enumerate}

\begin{lstlisting}
st = pd.read_csv("session_tracks.csv",
                 usecols=["session_id","user_id","position",
                          "track_id","playratio","label"])
st = st[st["track_id"].isin(i2v_ids)]
st = st[st["label"] != "unknown"]
st["playratio"] = st["playratio"].clip(upper=5.0)

keep = st.groupby("session_id")["track_id"].transform("count") >= 2
st = st[keep].sort_values(["session_id","position"]).reset_index(drop=True)
st.to_parquet("processed/session_tracks_i2v.parquet")
\end{lstlisting}

\textbf{Label semantics after cleaning.}
\begin{table}[H]
\centering
\caption{Label semantics and training weights.}
\begin{tabular}{llll}
\toprule
\textbf{Label} & \textbf{Typical playratio} & \textbf{Target $y$} & \textbf{Weight $\omega$} \\
\midrule
\texttt{positive} & $\ge 0.5$  & 1.0 & 1.0 \\
\texttt{neutral}  & 0.1--0.5   & 0.5 & 0.3 \\
\texttt{skip}     & $< 0.1$    & 0.0 & 1.0 \\
\texttt{unknown}  & null       & --- & Excluded \\
\bottomrule
\end{tabular}
\end{table}

\subsubsection{Remaining Interaction Tables}

\textbf{\texttt{session\_meta.csv}.}
Retain columns \texttt{session\_id} and \texttt{user\_id} only.
\texttt{session\_ts} is unreliable (not used for splitting; see Section~\ref{sec:split}).
Filter to sessions present in \texttt{session\_tracks\_i2v.parquet}.
Output: \texttt{processed/session\_meta\_i2v.parquet}.

\textbf{\texttt{playlist\_tracks.csv}.}
Retain \texttt{playlist\_id}, \texttt{user\_id}, \texttt{position}, \texttt{track\_id}.
Remove rows where \texttt{track\_id} $\notin$ \texttt{item2vec\_catalog}.
Drop playlists with remaining length $< 2$.
Output: \texttt{processed/playlist\_tracks\_i2v.parquet}.

\textbf{\texttt{playlist\_meta.csv}.}
Retain \texttt{playlist\_id} and \texttt{user\_id}.
Filter to playlists surviving the \texttt{playlist\_tracks} cleaning.
Output: \texttt{processed/playlist\_meta\_i2v.parquet}.

\textbf{\texttt{love.csv}.}
\begin{table}[H]
\centering
\caption{Column selection for \texttt{love.csv}.}
\begin{tabular}{lll}
\toprule
\textbf{Column} & \textbf{Keep?} & \textbf{Reason} \\
\midrule
\texttt{user\_id}  & Yes & Join key \\
\texttt{track\_id} & Yes & Explicit positive signal \\
\texttt{pref\_id}  & No  & Internal ID \\
\texttt{timestamp} & No  & Uniformly $-1$; no information content \\
\texttt{value}     & No  & Always ``love''; constant column \\
\bottomrule
\end{tabular}
\end{table}
Remove rows where \texttt{track\_id} $\notin$ \texttt{item2vec\_catalog}.
Remove duplicate \texttt{(user\_id, track\_id)} pairs.
Output: \texttt{processed/love\_filtered\_i2v.parquet}.

\textbf{\texttt{users.csv}.}
Retain \texttt{user\_id} only. All demographic columns have $\ge 25\%$ nulls
and are not used. Output: \texttt{processed/users\_i2v.parquet}.

\subsection{Expected Artifact Summary After All Preprocessing Stages}

\begin{table}[H]
\centering
\caption{Processed artifacts after Stages 1--7.}
\begin{tabular}{ll}
\toprule
\textbf{Artifact} & \textbf{Contents} \\
\midrule
\texttt{processed/catalog.csv} &
  Full track catalog with play counts and derived weights \\
\texttt{processed/item2vec\_catalog.csv} &
  Subset of catalog with Item2Vec embeddings  \\
\texttt{processed/item2vec\_128d.npy} &
  Embedding matrix, shape $(\text{vocab\_size}, 128)$ \\
\texttt{processed/item2vec\_track\_to\_row.json} &
  \texttt{track\_id} $\to$ row index mapping \\
\texttt{processed/item2vec\_model.bin} &
  Saved \texttt{KeyedVectors} for similarity queries \\
\texttt{processed/session\_tracks\_i2v.parquet} &
  Filtered session events (positive/neutral/skip, i2v tracks only) \\
\texttt{processed/session\_meta\_i2v.parquet} &
  Filtered session metadata \\
\texttt{processed/playlist\_tracks\_i2v.parquet} &
  Filtered playlist-track pairs \\
\texttt{processed/playlist\_meta\_i2v.parquet} &
  Filtered playlist metadata \\
\texttt{processed/love\_filtered\_i2v.parquet} &
  Filtered explicit love events \\
\texttt{processed/users\_i2v.parquet} &
  User IDs only \\
\texttt{processed/neg\_sample\_weights.npy} &
  Negative sampling weights for C3, length $=$ \texttt{vocab\_size} \\
\bottomrule
\end{tabular}
\end{table}

%% ===================================================================
\section{Data Split Strategy}
\label{sec:split}
%% ===================================================================

\subsection{Split Method}

\decision{Timestamp-based splitting is not used.
\texttt{love.csv} timestamps are uniformly $-1$.
\texttt{session\_ts} reliability in \texttt{session\_meta.csv} has not
been validated against the interaction data.
Sessions are split by \textbf{random session-level assignment}
in a 75\,/\,12.5\,/\,12.5 ratio with a fixed random seed.}

\begin{lstlisting}
import numpy as np, pandas as pd

session_ids = (pd.read_parquet("processed/session_meta_i2v.parquet")
               ["session_id"].values)

rng = np.random.default_rng(seed=42)
rng.shuffle(session_ids)

n = len(session_ids)
n_val  = int(0.125 * n)
n_test = int(0.125 * n)

test_sessions  = set(session_ids[:n_test])
val_sessions   = set(session_ids[n_test:n_test + n_val])
train_sessions = set(session_ids[n_test + n_val:])

# Save split assignments
np.save("processed/split_train.npy", np.array(list(train_sessions)))
np.save("processed/split_val.npy",   np.array(list(val_sessions)))
np.save("processed/split_test.npy",  np.array(list(test_sessions)))
\end{lstlisting}

\subsection{Usage per Split}

\begin{table}[H]
\centering
\caption{How each split is used across pipeline components.}
\begin{tabular}{lllr}
\toprule
\textbf{Split} & \textbf{Sessions} & \textbf{Usage} & \textbf{Est.\ contexts} \\
\midrule
Train & 75\% $\approx$ 1.69M &
  All prefix $\to$ next-track contexts for C3 training &
  $\approx$21M \\
Val & 12.5\% $\approx$ 281K &
  Last track of each session as evaluation label &
  $\approx$281K \\
Test & 12.5\% $\approx$ 281K &
  Last track for final evaluation &
  $\approx$281K \\
\bottomrule
\end{tabular}
\end{table}

\subsection{Evaluation Context Construction}

\paragraph{Training contexts.}
For each train session of length $N \ge 2$, iterate over positions
$t = 1, \ldots, N-2$, producing prefix $[i_1, \ldots, i_t]$ and
label $i_{t+1}$. Sessions of length 2 contribute exactly one context.
All three label types (positive, neutral, skip) in position $t+1$
contribute training rows with the weights defined in Table~5.

\paragraph{Validation and test contexts.}
For each val/test session of length $N \ge 2$:
\[
  \text{context} = [i_1, i_2, \ldots, i_{N-1}], \quad
  \text{label} = i_N.
\]

\note{No position-based hold-out (leave-last-out within a session) is
used. Sessions are split at the session level; no data leakage exists
across splits. The context for val/test is the full prefix up to the
penultimate track; the label is always the final track.}

\paragraph{Auxiliary data.}
Loved tracks (\texttt{love\_filtered\_i2v.parquet}) and playlist
memberships (\texttt{playlist\_tracks\_i2v.parquet}) are used for
K-Means user centroid computation regardless of session split,
as is standard practice in collaborative filtering evaluation.
Co-occurrence matrices are built from \textbf{training sessions only}.

%% ===================================================================
\section{Top-Level Architecture}
%% ===================================================================

The system recommends the next 5 tracks during an active listening session
on a Navidrome-based personal music server. The pipeline has four stages:

\begin{center}
\textbf{C1: Track Encoder (Item2Vec)}
$\;\to\;$
\textbf{C2: Multi-Recall Retriever}
$\;\to\;$
\textbf{C3: GRU Ranker}
$\;\to\;$
\textbf{C4: Policy Re-ranker}
\end{center}

C1 runs entirely offline; its outputs are materialised to disk and
frozen before any serving begins. C2 retrieval branches use precomputed
indices built from C1's outputs. C3 and C4 run at serving time.

\subsection{Latency Budget}

Production SLO: \textbf{p95 $< 300$\,ms} on a single CPU instance.
No GPU is required at serving time.

\begin{table}[H]
\centering
\caption{Latency budget allocation.}
\begin{tabular}{lrl}
\toprule
\textbf{Stage} & \textbf{Budget (ms)} & \textbf{Notes} \\
\midrule
Session context load (Redis) & 10 &
  User centroids + recent track embedding lookups \\
C2: Multi-Recall Retriever   & 80 &
  FAISS ANN + co-occurrence sparse lookup \\
C3: GRU Ranker               & 120 &
  GRU encode ($\approx$5\,ms) + 200$\times$MLP ($\approx$10\,ms) \\
C4: Policy Re-ranker         & 10  &
  Rule-based, no model calls \\
Response serialisation       & 10  & \\
\midrule
\textbf{Total} & \textbf{230} & 70\,ms headroom \\
\bottomrule
\end{tabular}
\end{table}

Hard timeout: \textbf{800\,ms} at the API gateway.

%% ===================================================================
\section{Component 1: Track Encoder}
%% ===================================================================

\subsection{Design Principle}

\decision{C1 has no learned parameters at serving time.
It is a pure offline preprocessing pipeline.
The only track feature is the 128-dimensional Item2Vec embedding.
Tag features and scalar metadata features are not used.
All outputs are materialised to disk and frozen before serving begins.}

C1's role is twofold:
\begin{enumerate}
  \item Produce a dense embedding $\mathbf{e}_i \in \mathbb{R}^{128}$
        for every track $i$ in the item2vec catalog.
        This embedding encodes the track's position in the
        session co-occurrence space.
  \item Build a FAISS approximate nearest-neighbour index over these
        embeddings for use by the Preference~NN retrieval branch (C2)
        and user centroid queries.
\end{enumerate}

The training of the embedding model is described in Stage~5 (B) of
the preprocessing pipeline. The description here focuses on the
serving-time contract.

\subsection{Embedding Lookup at Serving Time}

For a given \texttt{track\_id} $i$:
\begin{enumerate}
  \item Look up the row index $r = \texttt{track\_id\_to\_row}[i]$
        from \texttt{item2vec\_track\_to\_row.json}.
  \item Return $\mathbf{e}_i = \texttt{item2vec\_128d.npy}[r]$.
\end{enumerate}

Tracks not present in \texttt{item2vec\_track\_to\_row.json}
(i.e., tracks that appeared fewer than 5 times in all sessions)
have no embedding and are excluded from the system.

% \subsection{FAISS Index Construction}

% A FAISS index is built over all $(\text{vocab\_size}, 128)$
% Item2Vec embeddings. This index is used by:
% \begin{itemize}
%   \item The Preference NN retrieval branch (C2) at serving time.
%   \item User centroid computation (C2 offline) to verify centroid validity.
% \end{itemize}

% \begin{lstlisting}
% import faiss, numpy as np

% emb = np.load("processed/item2vec_128d.npy").astype("float32")
% d = emb.shape[1]   # 128

% index = faiss.index_factory(d, "IVF4096,PQ16")
% index.train(emb)
% index.add(emb)
% index.nprobe = 64

% faiss.write_index(index, "processed/faiss_item2vec.bin")
% print(f"Index built over {index.ntotal:,} vectors of dimension {d}")
% \end{lstlisting}

% \note{The IVF4096,PQ16 index type provides a good latency-recall trade-off
% at 128-d and catalog sizes of 400K--600K. \texttt{nprobe=64} (probing
% 64 of 4096 partitions) is set to balance recall and latency; tune if
% HR@200 falls below 0.60 on the validation set.}

\subsection{Summary of C1 Artifacts}

\begin{table}[H]
\centering
\caption{C1 artifacts and their consumers.}
\begin{tabular}{lll}
\toprule
\textbf{File} & \textbf{Shape / Type} & \textbf{Used by} \\
\midrule
\texttt{item2vec\_128d.npy}       & $(\text{N},\;128)$ float32 &
  C2 Pref NN, C2 centroids, C3 tokens, C3 scoring \\
\texttt{faiss\_item2vec.bin}      & FAISS index &
  C2 Pref NN ANN search \\
\texttt{item2vec\_track\_to\_row.json} & dict str $\to$ int &
  All components (embedding lookup) \\
\texttt{item2vec\_catalog.csv}    & tabular &
  Catalog filtering; metadata display \\
\bottomrule
\end{tabular}
\end{table}

%% ===================================================================
\section{Component 2: Multi-Recall Retriever}
%% ===================================================================

\subsection{Overview}

C2 produces a merged candidate set $\mathcal{C}_t$ with a hard cap of
200 candidates for the ranker. It operates through three independent
branches, each answering a different retrieval question:

\begin{table}[H]
\centering
\caption{Recall branches, signals, and candidate allocations.}
\begin{tabular}{llrl}
\toprule
\textbf{Branch} & \textbf{Signal} & \textbf{Candidates} &
  \textbf{Question answered} \\
\midrule
Co-occurrence & Session + playlist sequences & 100 &
  ``After track $X$, what typically follows?'' \\
Pref NN & Item2Vec similarity & 80 &
  ``Tracks with similar co-occurrence profile to your history'' \\
Popularity & Global event play count & 20 &
  ``Global safety net'' \\
\bottomrule
\end{tabular}
\end{table}

\note{ALS collaborative filtering was present in v3.0 but is excluded
here. The interaction matrix density of $0.0098\%$ makes ALS factors
unreliable for long-tail tracks. Item2Vec captures the collaborative
signal more robustly from the same interaction data.}

\subsection{Branch 1: Co-occurrence Retrieval (100 candidates)}

\subsubsection{Offline Construction}

\paragraph{Session transition matrix $C^{\text{sess}}$.}
From training sessions only (not val/test), for each adjacent pair
$(i_j, i_{j+1})$ where both have
\texttt{label} $\in \{\texttt{positive}, \texttt{neutral}\}$:
\[
  C^{\text{sess}}[i_j,\; i_{j+1}] \mathrel{+}= 1.
\]
Skip and unknown events are excluded to avoid reinforcing negative transitions.
Estimated: $\approx$21.9M adjacent positive/neutral pairs.
Storage format: SciPy sparse CSR (\texttt{.npz}), $\approx$200--500\,MB.

\paragraph{Playlist co-occurrence matrix $C^{\text{pl}}$.}
From all playlists (train and val/test, consistent with standard CF practice),
for each pair $(i, j)$ with $|\text{pos}_i - \text{pos}_j| \le 5$:
\[
  C^{\text{pl}}[i, j] \mathrel{+}= \frac{1}{|\text{pos}_i - \text{pos}_j|}.
\]
Playlist co-occurrence captures thematic proximity that session transitions
may miss (users curate stylistically coherent playlists).

\begin{lstlisting}
import scipy.sparse as sp, numpy as np, pandas as pd

st = pd.read_parquet("processed/session_tracks_i2v.parquet")
train_sessions = np.load("processed/split_train.npy", allow_pickle=True)
train_set = set(train_sessions)

st_train = st[st["session_id"].isin(train_set)].copy()
st_train = st_train[st_train["label"].isin({"positive","neutral"})]
st_train = st_train.sort_values(["session_id","position"])

# Build row/col index (track_id -> int)
track_ids = pd.read_csv("processed/item2vec_catalog.csv")["track_id"].values
tid2idx = {t: i for i, t in enumerate(track_ids)}
N = len(track_ids)

# Session transition matrix
rows, cols = [], []
for _, grp in st_train.groupby("session_id"):
    tids = grp["track_id"].tolist()
    for a, b in zip(tids[:-1], tids[1:]):
        if a in tid2idx and b in tid2idx:
            rows.append(tid2idx[a])
            cols.append(tid2idx[b])
data = np.ones(len(rows), dtype="float32")
C_sess = sp.csr_matrix((data, (rows, cols)), shape=(N, N))
sp.save_npz("processed/cooc_session.npz", C_sess)
\end{lstlisting}

\subsubsection{Serving}

Given the current session prefix $S_t = [i_{t-L+1}, \ldots, i_t]$
(up to $L = 10$ most recent tracks):
\begin{enumerate}
  \item Compute recency weight $w_j = \exp(-0.3 \cdot (t - j))$ for each track $i_j$.
  \item Score each candidate $c$:
        \[
          \text{score}_{\text{cooc}}(c) =
            \sum_{j \in S_t} w_j
            \bigl(
 C^{\text{sess}}[i_j, c]
              + C^{\text{pl}}[i_j, c]
            \bigr).
        \]
  \item Return top-100 by score.
\end{enumerate}
Serving latency: sparse row lookup is $O(\text{nnz per row})$, typically $< 2$\,ms.

\warn{If the current session contains only tracks absent from
$C^{\text{sess}}$ (cold tracks or long-tail), this branch may return
fewer than 100 candidates. The popularity branch compensates.}

\subsection{Branch 2: Preference Nearest Neighbours (80 candidates)}

\subsubsection{Multi-Centroid User Preference Vector}

A single mean embedding fails for users with diverse tastes
(e.g., classical and electronic music in the same history):
the mean falls in a region of the Item2Vec space that does not
correspond to any real track cluster. K-Means clustering preserves
the structure.

\paragraph{Offline construction (per user).}
\begin{enumerate}
  \item Collect Item2Vec embeddings for:
        \begin{itemize}
          \item All loved tracks (\texttt{love\_filtered\_i2v.parquet}),
                each counted $3\times$ to amplify explicit preference.
          \item All positive-label session tracks, capped at the
                most recent 300 events.
        \end{itemize}
        Tracks without Item2Vec embeddings are skipped.
  \item Determine $K$:
        $K = 1$ if the user has $< 10$ tracks;
        $K = 2$ if $10 \le \text{tracks} \le 30$;
        $K = 3$ if $> 30$ tracks.
  \item Run K-Means with chosen $K$ on the collected embeddings.
  \item Store per user: a list of up to 3 centroids
        $c_k \in \mathbb{R}^{128}$ with cluster sizes $n_k$.
\end{enumerate}

Cold users (no loved or positive tracks in the item2vec catalog):
no centroids are stored; this branch returns nothing at serving time.

\begin{lstlisting}
from sklearn.cluster import KMeans
import numpy as np, pandas as pd, pickle

emb = np.load("processed/item2vec_128d.npy")
t2r  = json.load(open("processed/item2vec_track_to_row.json"))

love = pd.read_parquet("processed/love_filtered_i2v.parquet")
sessions_pos = (pd.read_parquet("processed/session_tracks_i2v.parquet")
                .query("label == 'positive'")
                .sort_values(["user_id","session_id","position"])
                .groupby("user_id")
                .tail(300))

user_centroids = {}
for user_id in love["user_id"].unique():
    love_ids  = love.loc[love["user_id"] == user_id, "track_id"].tolist()
    pos_ids   = (sessions_pos[sessions_pos["user_id"] == user_id]
                 ["track_id"].tolist())
    track_ids = love_ids * 3 + pos_ids
    embs = np.array([emb[int(t2r[str(t)])]
                     for t in track_ids if str(t) in t2r])
    if len(embs) < 2:
        continue
    K = 1 if len(embs) < 10 else (2 if len(embs) <= 30 else 3)
    km = KMeans(n_clusters=K, random_state=42, n_init=10).fit(embs)
    user_centroids[user_id] = list(zip(km.cluster_centers_.tolist(),
                                       np.bincount(km.labels_).tolist()))

pickle.dump(user_centroids, open("processed/user_centroids.pkl","wb"))
\end{lstlisting}

\subsubsection{$\mathbf{u}^{\text{long}}$ Selection for C3}

At serving time, the Transformer receives a single long-term preference
vector. Select the centroid nearest to the current session's mean
Item2Vec embedding:
\[
  \mathbf{u}^{\text{long}}
  = \arg\min_{c_k}
    \bigl\| \bar{\mathbf{e}}^{\text{session}}_t - c_k \bigr\|_2,
\]
where $\bar{\mathbf{e}}^{\text{session}}_t$ is the mean of the current
session's track embeddings. This passes the taste dimension most
relevant to the current listening mood to C3's scoring head.

\subsubsection{Serving}

\begin{enumerate}
  \item Load the user's $K$ centroids from Redis.
  \item For each centroid $c_k$, query \texttt{faiss\_item2vec.bin}
        for $\lfloor 80 \cdot n_k / \sum n_k \rfloor$ nearest neighbours.
  \item Merge and deduplicate; return up to 80 candidates.
\end{enumerate}

FAISS ANN latency per centroid: $< 5$\,ms (IVF4096,PQ16 at nprobe=64).

\subsection{Branch 3: Popularity Fallback (20 candidates)}

Return the top-20 tracks by $\texttt{pop\_score} = \log(1 + \texttt{final\_playcount})$
that are not already in the candidate pool from the other branches.
This branch is precomputed as a sorted list at startup; serving latency
$< 1$\,ms. It ensures the candidate pool is never empty for cold users
with short sessions.

\subsection{Candidate Merge}

\begin{enumerate}
  \item Union of all branches: $100 + 80 + 20 = 200$ candidates before deduplication.
  \item Deduplicate by \texttt{track\_id}; if a track appears in multiple
        branches, keep once with the highest branch score.
  \item If $|\mathcal{C}_t| > 200$: retain top-200 by best rank across branches.
  \item If $|\mathcal{C}_t| < 50$: pad with additional popularity tracks
        (draws from the pre-sorted popularity list).
  \item Look up the 128-d Item2Vec embedding for each candidate from
        \texttt{item2vec\_128d.npy} (required by C3).
\end{enumerate}

%% ===================================================================
\section{Component 3: GRU Ranker}
%% ===================================================================

\subsection{Purpose}

C3 is the only component with \textbf{learned parameters}. It scores each
of the $\le 200$ candidates in $\mathcal{C}_t$ under the current session
context and user long-term preference, producing a ranked list from which
C4 selects the final Top-5.

\decision{The session encoder is a 2-layer GRU, not a Transformer.
The session median length is 6 tracks. At sequences of this length,
Transformer self-attention provides no meaningful advantage over GRU
recurrence, and the GRU eliminates the need for positional encodings.
GRU's sequential inductive bias is a natural fit for ordered listening
history, where the most recent track exerts the strongest influence.}

\subsection{Architecture}

\subsubsection{Token Construction}

For each track at position $j$ in the session prefix, the input token is:
\[
  \mathbf{x}_j
  = \mathbf{e}^{\text{i2v}}_{i_j} + \mathbf{e}^{\text{label}}_j
  \in \mathbb{R}^{128},
\]
where:
\begin{itemize}
  \item $\mathbf{e}^{\text{i2v}}_{i_j} \in \mathbb{R}^{128}$:
        the Item2Vec embedding of track $i_j$,
        looked up from \texttt{item2vec\_128d.npy}.
  \item $\mathbf{e}^{\text{label}}_j \in \mathbb{R}^{128}$:
        a \emph{learnable} label embedding;
        vocabulary $= \{\texttt{positive}, \texttt{neutral},
        \texttt{skip}, \texttt{pad}\}$.
        This allows C3 to weight skip events differently from
        positive ones in the session context.
\end{itemize}

Sessions shorter than the maximum length $L = 20$ are
\textbf{left-padded} with a learned \texttt{<PAD>} token
(assigned the \texttt{pad} label embedding).

\note{No positional encoding is needed. The GRU processes tokens
sequentially; position is encoded implicitly by the order of processing.
Earlier tracks influence the hidden state; later tracks override it.
The most recent track has the strongest influence on $\mathbf{s}_t$.}

\subsubsection{GRU Session Encoder}

\begin{table}[H]
\centering
\caption{GRU encoder hyperparameters.}
\begin{tabular}{lr}
\toprule
\textbf{Parameter} & \textbf{Value} \\
\midrule
Input dimension & 128 \\
Hidden dimension & 128 \\
Number of GRU layers & 2 \\
Dropout between layers & 0.1 \\
Batch-first & True \\
\bottomrule
\end{tabular}
\end{table}

\begin{lstlisting}
import torch, torch.nn as nn

class SessionEncoder(nn.Module):
    def __init__(self, d_emb=128, n_labels=4, dropout=0.1):
        super().__init__()
        self.label_emb = nn.Embedding(n_labels, d_emb)
        self.gru = nn.GRU(d_emb, d_emb,
                          num_layers=2,
                          batch_first=True,
                          dropout=dropout)

    def forward(self, item_embs, labels):
        # item_embs : (B, L, 128) -- Item2Vec embeddings
        # labels    : (B, L)     -- integer-coded, 0=positive...3=pad
        x = item_embs + self.label_emb(labels)  # (B, L, 128)
        _, h = self.gru(x)
        return h[-1]  # (B, 128) -- last GRU layer, last time step
\end{lstlisting}

The session representation is:
\[
  \mathbf{s}_t
  = \text{GRU}\!\bigl(\mathbf{x}_1, \mathbf{x}_2, \ldots, \mathbf{x}_L\bigr)
  \in \mathbb{R}^{128}.
\]

\subsubsection{Scoring Head}

For each candidate $i \in \mathcal{C}_t$:
\[
  \mathbf{z}_{t,i}
  = \bigl[
      \mathbf{s}_t \;;\;
      \mathbf{u}^{\text{long}} \;;\;
      \mathbf{e}^{\text{i2v}}_i \;;\;
      \bm{\phi}(t, i)
    \bigr]
  \in \mathbb{R}^{387}
\]

\begin{table}[H]
\centering
\caption{Scoring head input dimensions.}
\begin{tabular}{lr}
\toprule
\textbf{Component} & \textbf{Dim} \\
\midrule
Session representation $\mathbf{s}_t$ & 128 \\
User long-term preference $\mathbf{u}^{\text{long}}$ (nearest centroid) & 128 \\
Candidate Item2Vec embedding $\mathbf{e}^{\text{i2v}}_i$ & 128 \\
Cross features $\bm{\phi}(t, i)$ & 3 \\
\midrule
\textbf{Total} & \textbf{387} \\
\bottomrule
\end{tabular}
\end{table}

\paragraph{Cross features $\bm{\phi}(t, i)$ (3-d).}
These give the MLP direct access to the most discriminative similarity
signals, without requiring the MLP to discover them implicitly.
All similarities are computed in the Item2Vec embedding space.
\begin{enumerate}
  \item $\cos\!\bigl(\mathbf{e}^{\text{i2v}}_i,\;
        \mathbf{e}^{\text{i2v}}_{i_t}\bigr)$:
        co-occurrence similarity with the most recently played track.
  \item $\cos\!\bigl(\mathbf{e}^{\text{i2v}}_i,\;
        \bar{\mathbf{e}}^{\text{session}}_t\bigr)$:
        co-occurrence similarity with the session mean embedding.
  \item $\cos\!\bigl(\mathbf{e}^{\text{i2v}}_i,\;
        \mathbf{u}^{\text{long}}\bigr)$:
        co-occurrence similarity with the user's selected long-term centroid.
\end{enumerate}

\paragraph{MLP layers.}
\[
  387
  \;\xrightarrow{\text{Linear, ReLU, drop 0.1}}\;
  256
  \;\xrightarrow{\text{Linear, ReLU, drop 0.1}}\;
  64
  \;\xrightarrow{\text{Linear}}\;
  1
\]

The scalar output $\hat{y}_{t,i}$ is the predicted relevance of
candidate $i$ as the next track given session context $t$.

\subsection{Training}
\label{sec:c3training}

\subsubsection{Key Design Decision: Full-Catalog Negative Sampling}

\decision{C3 training data is constructed by sampling from the full
item2vec catalog, not from C2's retrieval output.}

If training were restricted to C2's 200 candidates, C3 could only
learn to re-rank what C2 already retrieved. Tracks outside C2's
recall ceiling would never appear as positives during training,
hard-capping C3's generalisation. By sampling from the full catalog,
C3 learns a general relevance function that transfers to any set of
candidates presented by C2 at serving time.

\subsubsection{Training Sample Construction}

For each training session of length $N \ge 2$, iterate over positions
$t = 1, \ldots, N-2$:

\begin{enumerate}
  \item \textbf{Prefix:} $[i_1, \ldots, i_t]$, left-padded to length $L = 20$.
  \item \textbf{Positive:} $i_{t+1}$ (the actual next track).
        \begin{itemize}
          \item \texttt{positive}: $y = 1.0$, $\omega = 1.0$
          \item \texttt{neutral}: $y = 0.5$, $\omega = 0.3$
          \item \texttt{skip}: $y = 0.0$, $\omega = 1.0$
          \item \texttt{unknown}: skip this context entirely.
        \end{itemize}
  \item \textbf{Hard negatives (3):}
        sampled from $C^{\text{sess}}$'s top-neighbours of recent session
        tracks, excluding $i_{t+1}$. These are plausible-but-wrong candidates
        that challenge the ranker to discriminate.
  \item \textbf{Random negatives (2):}
        sampled from the full item2vec catalog
        $\propto \texttt{neg\_sample\_weight} = \texttt{final\_playcount}^{0.75}$,
        normalised. Popularity-weighted sampling ensures negatives are
        not trivially easy (avoiding extreme long-tail items).
\end{enumerate}

Each context produces 6 rows (1 positive + 3 hard + 2 random negatives).
Estimated total training rows: $\approx$21M contexts $\times$ 6 = $\approx$126M rows.

\subsubsection{Training Objective}

Joint pointwise and pairwise loss:
\[
  \mathcal{L} = \mathcal{L}_{\text{point}} + 0.5 \cdot \mathcal{L}_{\text{pair}}.
\]

\paragraph{Pointwise (weighted BCE):}
\[
  \mathcal{L}_{\text{point}}
  = \sum_{(t,i)} \omega_{t,i} \cdot
    \mathrm{BCE}\!\bigl(\hat{y}_{t,i},\; y_{t,i}\bigr).
\]

\paragraph{Pairwise (BPR):}
For each (positive, negative) pair $(i^+, i^-)$ in the same context:
\[
  \mathcal{L}_{\text{pair}}
  = -\log \sigma\!\bigl(
      \hat{y}_{t,i^+} - \hat{y}_{t,i^-}
    \bigr).
\]

BPR directly optimises the ranking order between the positive and each
negative, complementing the pointwise calibration signal from BCE.

\subsubsection{Training Hyperparameters}

\begin{table}[H]
\centering
\caption{GRU ranker training hyperparameters.}
\begin{tabular}{lr}
\toprule
\textbf{Hyperparameter} & \textbf{Value} \\
\midrule
Batch size & 512 contexts ($\times 6 = 3072$ rows) \\
Optimizer & AdamW \\
Learning rate & $1 \times 10^{-4}$ \\
Weight decay & $10^{-4}$ \\
LR schedule & Cosine annealing, 10\% linear warmup \\
Epochs & 3 \\
Gradient clipping & max-norm $= 1.0$ \\
Max sequence length $L$ & 20 \\
\bottomrule
\end{tabular}
\end{table}

\subsubsection{Inference (Serving)}

\begin{enumerate}
  \item Look up Item2Vec embeddings for all tracks in the session prefix.
        Construct token inputs $\{\mathbf{x}_j\}$ with label embeddings.
  \item Forward pass through the 2-layer GRU $\to \mathbf{s}_t$
        ($\approx$5\,ms).
  \item Retrieve $\mathbf{u}^{\text{long}}$ from Redis.
  \item For all $|\mathcal{C}_t| \le 200$ candidates:
        look up embeddings, compute 3 cross features,
        batch into one MLP forward pass ($\approx$10\,ms).
  \item Sort candidates by $\hat{y}_{t,i}$.
\end{enumerate}

Total C3 latency: $\approx$15--20\,ms, within the 120\,ms budget.

% %% ===================================================================
% \section{Component 4: Policy Re-ranker}
% %% ===================================================================

% \subsection{Policy Rules}

% C4 has no learned parameters. It applies deterministic rule-based
% transformations to C3's ranked output to enforce business constraints
% and inject diversity.

% \paragraph{Step 1: Hard filters.}
% Remove any candidate that:
% \begin{itemize}
%   \item Was played in the last 30 tracks of this session.
%   \item Has been explicitly disliked by the user.
% \end{itemize}
% Applied via a penalty of $-\infty$ (score $= -100.0$) to effectively
% eliminate these tracks from selection.

% \paragraph{Step 2: Same-artist soft penalty.}
% \[
%   s'_i = s_i - 0.15 \cdot \mathbb{1}[\text{artist}(i) \in
%     \{\text{artist}(i_{t-1}), \text{artist}(i_t)\}].
% \]
% Avoids back-to-back same-artist recommendations without hard-blocking
% same-artist tracks from appearing later in the queue.

% \paragraph{Step 3: Greedy MMR diversity.}
% Select Top-5 one track at a time. After each selection, apply a
% diversity bonus to remaining candidates:
% \[
%   s''_i = s'_i + \delta \cdot \Bigl(1 - \max_{j \in \text{selected}}
%     \cos\!\bigl(\mathbf{e}^{\text{i2v}}_i,\;
%                 \mathbf{e}^{\text{i2v}}_j\bigr)\Bigr),
%   \quad \delta = 0.05.
% \]
% Candidates less similar to already-selected tracks receive a small bonus,
% promoting co-occurrence-space diversity in the final list.

% \paragraph{Step 4: $\epsilon$-greedy exploration.}
% The 5th slot (last position) is replaced with a uniformly random
% candidate from $\mathcal{C}_t \setminus \{\text{Top-4}\}$ with
% probability $\epsilon = 0.10$. This provides exploration signal
% for online retraining.

% \begin{table}[H]
% \centering
% \caption{Policy parameters.}
% \begin{tabular}{llll}
% \toprule
% \textbf{Parameter} & \textbf{Symbol} & \textbf{Value} & \textbf{Tuning} \\
% \midrule
% No-repeat penalty & $\alpha$ & $-100.0$ & Fixed (business rule) \\
% Same-artist penalty & $\beta$ & 0.15 & Grid search on val skip-rate \\
% Diversity bonus & $\delta$ & 0.05 & Manual, A/B experiment \\
% Exploration probability & $\epsilon$ & 0.10 & Fixed \\
% Recency window & $M$ & 30 tracks & Fixed \\
% \bottomrule
% \end{tabular}
% \end{table}

% \subsection{Fallback Chain}

% \begin{enumerate}
%   \item \textbf{Primary:} Full pipeline (C2 $\to$ C3 $\to$ C4).
%   \item \textbf{Fallback 1:} C3 timeout
%         $\to$ return C2's top-5 by co-occurrence score.
%   \item \textbf{Fallback 2:} C2 timeout
%         $\to$ top-10 Item2Vec nearest neighbours of the most recent track
%         (precomputed, cached in Redis).
%   \item \textbf{Fallback 3:} All else fails
%         $\to$ top-5 by global \texttt{pop\_score} (precomputed sorted list).
% \end{enumerate}

% %% ===================================================================
% \section{Full Training Pipeline}
% %% ===================================================================

% The end-to-end training pipeline has nine phases (Phase~0 through Phase~8).
% All phases are CPU-only except Phase~6 (GRU training), which requires a GPU.

% \subsection{Phase 0: Data Preparation (CPU, 2--4 h)}

% \begin{enumerate}
%   \item Run Stage~1: deduplicate \texttt{tracks.csv}
%         $\to$ \texttt{tracks\_unique.csv}.
%   \item Run Stage~2: aggregate \texttt{events.csv}
%         $\to$ \texttt{event\_playcount.csv}.
%   \item Run Stage~3: merge into \texttt{catalog.csv} with
%         \texttt{pop\_score} and \texttt{neg\_sample\_weight}.
% \end{enumerate}

% \textbf{Checkpoint:} Print catalog size and the fraction of tracks
% with \texttt{final\_playcount} $= 0$.

% \subsection{Phase 1: Item2Vec Training (CPU, $\approx$1 h)}

% \begin{enumerate}
%   \item Run Stage~4 (A): build \texttt{item2vec\_corpus.parquet}.
%         Log session count, token count, and raw vocabulary size.
%   \item Run Stage~5 (B): train Word2Vec Skip-gram model.
%         Log vocabulary size after \texttt{min\_count=5}, training duration.
%   \item Save \texttt{item2vec\_128d.npy}, \texttt{item2vec\_track\_to\_row.json},
%         \texttt{item2vec\_catalog.csv}.
%   \item Build FAISS index \texttt{faiss\_item2vec.bin}.
% \end{enumerate}

% \textbf{Checkpoint:} Run Stage~6 (C) validation.
% Verify: no NaN in embeddings; same-artist similarity $>$ random;
% at least one nearest-neighbour per sampled track is from the same artist.

% \subsection{Phase 2: Interaction Filtering (CPU, $\approx$30 min)}

% Run Stage~7 (D): filter all interaction tables to item2vec catalog.
% Output the \texttt{\_i2v.parquet} files.

% \textbf{Checkpoint:} Print per-table row counts before and after filtering.
% Print the fraction of session events covered by the item2vec catalog.

% \subsection{Phase 3: Data Split (CPU, $< 5$ min)}

% Random session-level split: 75 / 12.5 / 12.5.
% Save split sets to \texttt{split\_train.npy}, \texttt{split\_val.npy},
% \texttt{split\_test.npy}.

% \textbf{Checkpoint:} Verify no session\_id appears in two splits.
% Print session and interaction counts per split.

% \subsection{Phase 4: Co-occurrence Matrix Construction (CPU, $\approx$1 h)}

% Build $C^{\text{sess}}$ (from training sessions, positive + neutral pairs)
% and $C^{\text{pl}}$ (from all playlists, weighted by position distance).
% Save both as SciPy CSR \texttt{.npz}.

% \textbf{Checkpoint:} Print matrix densities and non-zero counts.
% Verify that a known same-session pair has a non-zero entry.

% \subsection{Phase 5: User Centroid Computation (CPU, $\approx$15 min)}

% For each user, collect positive/loved track embeddings, run K-Means
% ($K \le 3$), and save centroids to \texttt{user\_centroids.pkl}.
% Load into Redis as \texttt{user:\{user\_id\}:centroids}.

% \textbf{Checkpoint:} Print the fraction of users with at least one centroid.
% Verify K-Means convergence warning rate is low.

% \subsection{Phase 6: Ranker Training Data Generation (CPU, 6--10 h)}

% For each training session prefix:
% \begin{enumerate}
%   \item Run the full C2 pipeline (co-occurrence + Pref~NN + popularity)
%         to confirm hard-negative pool availability.
%   \item Sample 3 hard negatives from $C^{\text{sess}}$ neighbours.
%   \item Sample 2 random negatives $\propto \texttt{neg\_sample\_weight}$.
%   \item Attach labels, weights, cross-features.
% \end{enumerate}
% Save \texttt{ranker\_train.parquet} and \texttt{ranker\_val.parquet}.

% \textbf{Checkpoint:} Verify positive track matches actual next track.
% No negative has the same \texttt{track\_id} as the positive.
% Print label distribution across all rows.

% \subsection{Phase 7: GRU Ranker Training (GPU, 4--8 h)}

% \begin{enumerate}
%   \item Initialise GRU encoder and scoring MLP (Xavier initialisation).
%   \item Train 3 epochs with joint BCE~+~BPR loss.
%   \item Evaluate on \texttt{ranker\_val.parquet} after each epoch:
%         HR@5, NDCG@5, MRR@5.
%   \item Save best checkpoint by NDCG@5.
% \end{enumerate}

% \textbf{Checkpoint:} NDCG@5 $\ge 0.12$ on validation set.
% HR@5 $\ge 0.20$.
% Ranker must improve over retriever-only (co-occurrence top-5) baseline.

% \subsection{Phase 8: End-to-End Evaluation and Deployment (CPU, 4--6 h)}

% \begin{enumerate}
%   \item Run the full pipeline on the test split.
%         Report all metrics (Table~\ref{tab:eval_metrics}) separately
%         for warm ($\ge 50$ train events), light (1--49), and cold (0) users.
%   \item Latency profiling: 1000 random serving requests on a single CPU core.
%   \item Ablation experiments:
%         \begin{itemize}
%           \item Ranker vs.\ retriever-only baseline.
%           \item With vs.\ without co-occurrence branch.
%           \item With vs.\ without Pref~NN branch.
%           \item With vs.\ without MMR diversity.
%           \item Item2Vec vs.\ popularity-only baseline (lower bound).
%         \end{itemize}
%   \item Package all artefacts into a versioned directory:
% \begin{lstlisting}
% artifacts/v1.0.0/
%   item2vec_128d.npy
%   item2vec_track_to_row.json
%   faiss_item2vec.bin
%   item2vec_catalog.csv
%   cooc_session.npz
%   cooc_playlist.npz
%   user_centroids.pkl
%   gru_ranker.pt
%   gru_ranker_config.json
%   neg_sample_weights.npy
%   metadata.json           # version, date, metrics, split sizes
% \end{lstlisting}
% \end{enumerate}

% \textbf{Go / no-go criteria:}
% HR@5 $\ge 0.25$; p95 latency $\le 300$\,ms; diversity $\ge 0.40$.

% %% ===================================================================
% \section{Logging Schema for Online Retraining}
% %% ===================================================================

% \subsection{Impression Log}

% Written to Postgres on every recommendation call.

% \begin{table}[H]
% \centering
% \caption{Impression log schema.}
% \begin{tabular}{lll}
% \toprule
% \textbf{Column} & \textbf{Type} & \textbf{Description} \\
% \midrule
% \texttt{request\_id}          & UUID      & \\
% \texttt{user\_id}             & INT       & \\
% \texttt{session\_id}          & VARCHAR   & \\
% \texttt{timestamp}            & TIMESTAMP & \\
% \texttt{trigger\_type}        & ENUM      & \texttt{END} or \texttt{SKIP} \\
% \texttt{context\_track\_ids}  & INT[]     & Recent tracks (up to 20) \\
% \texttt{candidate\_pool\_ids} & INT[]     & All $\le$200 candidates \\
% \texttt{candidate\_scores}    & FLOAT[]   & GRU ranker scores \\
% \texttt{top5\_ids}            & INT[]     & Final returned IDs \\
% \texttt{top5\_scores}         & FLOAT[]   & \\
% \texttt{exploration\_pos}     & INT[]     & Positions filled by $\epsilon$-greedy \\
% \texttt{model\_version}       & VARCHAR   & \\
% \texttt{fallback\_level}      & INT       & 0 (primary) -- 3 (global pop) \\
% \texttt{latency\_ms}          & FLOAT     & \\
% \bottomrule
% \end{tabular}
% \end{table}

% \subsection{Outcome Log}

% Written when the next playback event occurs after a recommendation.

% \begin{table}[H]
% \centering
% \caption{Outcome log schema.}
% \begin{tabular}{lll}
% \toprule
% \textbf{Column} & \textbf{Type} & \textbf{Description} \\
% \midrule
% \texttt{outcome\_id}       & UUID   & \\
% \texttt{request\_id}       & UUID   & FK $\to$ impression log \\
% \texttt{chosen\_track\_id} & INT    & What the user played next \\
% \texttt{chosen\_from\_rec} & BOOL   & Whether it was in the Top-5 \\
% \texttt{chosen\_position}  & INT    & Position in Top-5; $-1$ if not \\
% \texttt{play\_duration\_sec} & FLOAT & \\
% \texttt{play\_ratio}       & FLOAT  & \\
% \texttt{explicit\_feedback} & ENUM  & \texttt{like}, \texttt{dislike}, \texttt{none} \\
% \texttt{derived\_label}    & ENUM   & \texttt{positive}, \texttt{skip}, \texttt{neutral}, \texttt{unknown} \\
% \bottomrule
% \end{tabular}
% \end{table}

% \subsection{Retraining Trigger and Promotion}

% \paragraph{Retraining triggers.}
% \begin{itemize}
%   \item Scheduled: weekly.
%   \item Triggered: 7-day rolling skip rate increases $> 3$ percentage points;
%         fallback rate $> 10\%$; or average play ratio on recommended tracks
%         drops below 0.60.
% \end{itemize}

% \paragraph{Two-gate promotion.}
% \begin{enumerate}
%   \item \textbf{Offline gate:} New model achieves $\ge 95\%$ of the current
%         model's NDCG@5 on a 1-week holdout.
%   \item \textbf{Shadow gate:} 24-hour shadow scoring;
%         promote if shadow scores correlate with positive outcomes at least
%         as well as production scores.
% \end{enumerate}

% \paragraph{Rollback.}
% If the promoted model's skip rate increases $> 5$ percentage points
% within 48 hours, auto-revert to the previous checkpoint.

% %% ===================================================================
% \section{Evaluation Metrics}
% %% ===================================================================

% \begin{table}[H]
% \centering
% \caption{Offline and online evaluation metrics.}
% \label{tab:eval_metrics}
% \begin{tabular}{llll}
% \toprule
% \textbf{Metric} & \textbf{Scope} & \textbf{Target} & \textbf{Notes} \\
% \midrule
% HR@5          & Retriever + Ranker & $\ge 0.25$ &
%   Fraction of test contexts where the label is in Top-5 \\
% NDCG@5        & Ranker output      & $\ge 0.15$ &
%   Normalised discounted cumulative gain \\
% MRR@5         & Ranker output      & $\ge 0.15$ &
%   Mean reciprocal rank \\
% HR@200        & C2 Retriever       & $\ge 0.60$ &
%   Recall of the label in the 200 candidates \\
% Skip rate     & Online             & $\le 15\%$ &
%   Fraction of recommended tracks played $< 5$\,s \\
% Play ratio    & Online             & $\ge 0.60$ &
%   Mean ratio of track duration played \\
% Diversity     & Policy output      & $\ge 0.40$ &
%   $1 - $ mean pairwise cosine similarity of Top-5 in Item2Vec space \\
% Fallback rate & Serving            & $\le 5\%$  &
%   Fraction of requests that fall back beyond primary pipeline \\
% p95 latency   & Serving            & $\le 300$\,ms & Single CPU instance \\
% \bottomrule
% \end{tabular}
% \end{table}

% Report HR@5, NDCG@5, MRR@5 separately for:
% \begin{itemize}
%   \item \textbf{Warm} users: $\ge 50$ interaction events in training data.
%   \item \textbf{Light} users: 1--49 events.
%   \item \textbf{Cold} users: 0 events (no training history).
% \end{itemize}

% %% ===================================================================
% \section{Summary of Key Decisions}
% %% ===================================================================

% \begin{table}[H]
% \centering
% \caption{Index of all binding architectural decisions.}
% \begin{tabular}{lp{11cm}}
% \toprule
% \textbf{ID} & \textbf{Decision} \\
% \midrule
% D1 & Item2Vec (Skip-gram, 128-d) replaces PANNs audio embeddings.
%      Justification: PANNs requires $\sim$10 days of downloads with 30\%
%      hit rate; Item2Vec trains in $\sim$1 hour on local data. \\
% D2 & Item2Vec is trained on \textbf{all available sessions}
%      (not restricted to the training split), because it is unsupervised
%      and carries no label-leakage risk. \\
% D3 & Item2Vec catalog is defined by \texttt{min\_count=5}: only tracks
%      with $\ge 5$ session appearances receive an embedding. \\
% D4 & ALS collaborative filtering is excluded. Interaction matrix density
%      of 0.0098\% makes ALS factors unreliable for long-tail tracks. \\
% D5 & Tag features are excluded. 70.5\% of tracks have no tag annotations. \\
% D6 & Scalar metadata features are excluded. The only track representation
%      is the 128-d Item2Vec embedding. \\
% D7 & Three recall branches: co-occurrence (100) + Pref~NN (80) +
%      popularity (20). No ALS branch. \\
% D8 & Pref~NN FAISS index is built over Item2Vec embeddings (not audio).
%      Pref~NN therefore captures co-occurrence-space similarity, not
%      acoustic similarity. \\
% D9 & Multi-centroid K-Means ($K \le 3$) for user long-term preference.
%      Single mean embedding loses cluster structure for diverse users. \\
% D10 & $\mathbf{u}^{\text{long}}$ at serving time is the centroid nearest
%       to the current session's mean embedding, not a fixed centroid. \\
% D11 & Session encoder is a 2-layer GRU, not a Transformer.
%       Session median length is 6; GRU provides equivalent quality
%       with simpler serving code and no positional encoding. \\
% D12 & C3 training uses full-catalog negative sampling to avoid
%       hard-capping on C2's recall ceiling. \\
% D13 & Joint loss: BCE (pointwise, weighted by label) + 0.5$\times$BPR (pairwise). \\
% D14 & Data split is random session-level (75/12.5/12.5).
%       Timestamps are unreliable; temporal split is not used. \\
% D15 & C4: hard-block replays and dislikes;
%       same-artist soft penalty ($\beta = 0.15$);
%       greedy MMR diversity ($\delta = 0.05$);
%       $\epsilon$-greedy exploration ($\epsilon = 0.10$) at position~5. \\
% D16 & Retraining is weekly by default, with skip-rate and play-ratio
%       KPI triggers. Two-gate promotion with auto-rollback. \\
% \bottomrule
% \end{tabular}
% \end{table}

\end{document}